\documentclass{article}
\usepackage{amsmath,amssymb,amsthm}
\usepackage{algorithm}
\usepackage{algpseudocode}
\usepackage{geometry}
\geometry{margin=1in}
\newtheorem{theorem}{Theorem}
\newtheorem{lemma}{Lemma}
\newcommand{\E}{\mathbb{E}}
\newcommand{\R}{\mathbb{R}}
\newcommand{\norm}[1]{\left\lVert #1 \right\rVert}
\begin{document}

\section*{AdaGrad (Adaptive Gradient) }

\section*{Mathematical Formulation}
Let $f:\R^{d}\rightarrow\R$ be a convex and $L$-smooth objective.  
Denote by $g_{k}=\nabla f(w_{k})$ the (full or stochastic) gradient at iteration $k$.  
Define the element‑wise accumulated squared gradients
\[
G_{k}\;:=\;\sum_{i=1}^{k} g_{i}\odot g_{i}\in\R^{d},
\qquad 
\bigl(G_{k}\bigr)_{j}=\sum_{i=1}^{k} g_{i,j}^{2},
\]
where $\odot$ denotes the Hadamard (element‑wise) product.  
For a positive stepsize parameter $\alpha>0$ and a small constant $\varepsilon>0$ the AdaGrad update is
\[
\boxed{
w_{k+1}= w_{k}\;-\;\alpha\;\frac{g_{k}}{\sqrt{G_{k}+\varepsilon}}\;,
}
\tag{1}
\]
where the division is also element‑wise:
\[
\bigl(\alpha/\sqrt{G_{k}+\varepsilon}\bigr)_{j}
   =\frac{\alpha}{\sqrt{(G_{k})_{j}+\varepsilon}} .
\]

\section*{Pseudocode}
\begin{algorithm}[H]
\caption{AdaGrad}
\begin{algorithmic}[1]
\Require Initial point $w_{0}\in\R^{d}$, stepsize $\alpha>0$, tolerance $\varepsilon>0$
\State $G_{0}\gets 0\in\R^{d}$ \Comment{accumulated squared gradients}
\For{$k=0,1,\dots,T-1$}
    \State Compute (stochastic) gradient $g_{k}= \nabla f(w_{k})$
    \State $G_{k+1}\gets G_{k}+g_{k}\odot g_{k}$ \Comment{update accumulator}
    \State $ \eta_{k}\gets \alpha\;/\;\sqrt{G_{k+1}+\varepsilon}$ \Comment{element‑wise stepsize}
    \State $w_{k+1}\gets w_{k}-\eta_{k}\odot g_{k}$ \Comment{parameter update}
\EndFor
\State \Return $w_{T}$
\end{algorithmic}
\end{algorithm}

\section*{Dry Run Example}
Consider the one‑dimensional quadratic $f(w)=(w-3)^{2}$, $w_{0}=0$, $\alpha=0.1$, $\varepsilon=10^{-8}$.

\begin{center}
\begin{tabular}{c|c|c|c|c}
$k$ & $w_{k}$ & $g_{k}=f'(w_{k})$ & $G_{k}$ (accumulated squares) & $f(w_{k})$\\ \hline
0 & $0$ & $2(0-3)=-6$ & $0$ & $9$\\
1 & $0-0.1\frac{-6}{\sqrt{36+\varepsilon}}=0+0.1$ & $2(0.1-3)=-5.8$ & $36$ & $8.41$\\
2 & $0.1-0.1\frac{-5.8}{\sqrt{36+33.64+\varepsilon}}=0.1+0.1\frac{5.8}{\sqrt{69.64}}\approx0.1699$ &
 $2(0.1699-3)=-5.6602$ &
 $36+33.64=69.64$ &
 $8.025$\\
3 & $0.1699-0.1\frac{-5.6602}{\sqrt{69.64+32.048}} 
      \approx0.1699+0.1\frac{5.6602}{\sqrt{101.688}}
      \approx0.2270$ &
 $2(0.2270-3)=-5.5460$ &
 $69.64+32.048\approx101.688$ &
 $7.613$\\
\end{tabular}
\end{center}
The objective value decreases from $9$ to $7.613$ after three iterations, illustrating the adaptive reduction of the stepsize as the accumulated squared gradients grow.

\newpage
\section*{Assumptions}
\begin{enumerate}
\item \textbf{Convexity.} $f$ is convex:
\[
f(u)\ge f(v)+\langle\nabla f(v),u-v\rangle,\qquad\forall u,v\in\R^{d}.
\tag{A1}
\]
\item \textbf{L‑smoothness.} $\nabla f$ is $L$‑Lipschitz:
\[
\norm{\nabla f(u)-\nabla f(v)}\le L\norm{u-v},\qquad\forall u,v\in\R^{d}.
\tag{A2}
\]
\item \textbf{Bounded gradients.} For all $k$, $\norm{g_{k}}\le G_{\max}$.
\item \textbf{Existence of optimal solution.} There exists $w^{*}\in\arg\min_{w}f(w)$ and its distance to the start point is bounded:
\[
D:=\norm{w_{0}-w^{*}}<\infty .
\tag{A3}
\]
\item \textbf{Non‑negative accumulator.} $\varepsilon>0$ guarantees $\sqrt{G_{k}+\varepsilon}$ is well defined.
\end{enumerate}

\section*{Main Convergence Theorem}
\begin{theorem}[AdaGrad Regret / Suboptimality Bound]\label{thm:adagrad}
Under Assumptions (A1)–(A4) and using the update \eqref{1} with $\alpha>0$, the iterates satisfy for any $T\ge 1$
\[
\frac{1}{T}\sum_{k=0}^{T-1}\bigl(f(w_{k})-f(w^{*})\bigr)
   \;\le\;
   \frac{D^{2}}{2\alpha T}\sum_{j=1}^{d}\sqrt{G_{T,j}+\varepsilon}
   \;+\;
   \frac{\alpha G_{\max}^{2}}{2T}\sum_{j=1}^{d}\frac{1}{\sqrt{G_{T,j}+\varepsilon}} .
\tag{2}
\]
In particular, because $\sqrt{G_{T,j}+\varepsilon}\le\sqrt{T}\,G_{\max}$, we obtain the $O\!\bigl(\frac{1}{\sqrt{T}}\bigr)$ rate
\[
\frac{1}{T}\sum_{k=0}^{T-1}\bigl(f(w_{k})-f(w^{*})\bigr)
   \;\le\;
   \frac{D G_{\max}\sqrt{d}}{\sqrt{T}} .
\tag{3}
\]
\end{theorem}

\section*{Theoretical Properties}
\begin{itemize}
\item \textbf{Stability.} The denominator $\sqrt{G_{k}+\varepsilon}$ never vanishes, thus the stepsize is bounded above by $\alpha/\sqrt{\varepsilon}$ and shrinks monotonically.
\item \textbf{Hyperparameter Sensitivity.} The only free scalar is $\alpha$; the bound \eqref{2} is linear in $\alpha$ for the second term and inversely proportional for the first term, suggesting the optimal $\alpha^{*}=D G_{\max}\sqrt{d}$ up to constants.
\item \textbf{Comparison to SGD.} For a fixed stepsize $\eta$, SGD yields $O(1/\sqrt{T})$ in expectation but with a constant depending on $\eta$. AdaGrad automatically adapts to the geometry of the data, achieving the same asymptotic rate without manual tuning of a decaying schedule.
\end{itemize}

\section*{Discussion}
The theorem shows that AdaGrad attains the optimal $O(1/\sqrt{T})$ convergence for convex smooth problems while being fully adaptive. The bound depends on the accumulated squared gradients per coordinate; if some coordinates receive very small gradients, their stepsizes remain relatively large, which explains the empirical advantage on sparse data.

\newpage
\section*{Preliminaries (Definitions and Lemmas)}
\begin{lemma}[Three–point identity for Euclidean norm]\label{lem:threepoint}
For any $a,b,c\in\R^{d}$ and any positive diagonal matrix $H\succ0$,
\[
\|a-c\|_{H}^{2}
   =\|a-b\|_{H}^{2}+\|b-c\|_{H}^{2}+2\langle a-b, H(b-c)\rangle .
\tag{4}
\]
\end{lemma}
\begin{proof}
Expand the squared norm $\|x\|_{H}^{2}=x^{\top}Hx$ and collect terms. \hfill $\square$
\end{proof}

\section*{Main Proof (Step‑by‑Step Derivation)}
We prove Theorem~\ref{thm:adagrad}. Throughout the proof we denote the diagonal matrix
\[
H_{k}:=\operatorname{diag}\!\bigl(\sqrt{G_{k}+\varepsilon}\bigr) .
\]

\noindent\textbf{[Step 1]}  Start from the convexity inequality (A1) with $u=w^{*}$ and $v=w_{k}$:
\[
f(w_{k})-f(w^{*})\le\langle g_{k},\,w_{k}-w^{*}\rangle .
\tag{5}
\]

\noindent\textbf{[Step 2]}  Rewrite the inner product using the update \eqref{1}. Since $w_{k+1}=w_{k}-\alpha H_{k}^{-1}g_{k}$,
\[
w_{k}-w^{*}=w_{k+1}-w^{*}+\alpha H_{k}^{-1}g_{k}.
\tag{6}
\]

\noindent\textbf{[Step 3]}  Substitute \eqref{6} into \eqref{5}:
\[
f(w_{k})-f(w^{*})
   \le\langle g_{k},\,w_{k+1}-w^{*}\rangle
      +\alpha\langle g_{k},\,H_{k}^{-1}g_{k}\rangle .
\tag{7}
\]

\noindent\textbf{[Step 4]}  Apply Lemma~\ref{lem:threepoint} with $a=w_{k}$, $b=w_{k+1}$, $c=w^{*}$ and $H=H_{k}$:
\[
\|w_{k}-w^{*}\|_{H_{k}}^{2}
   =\|w_{k+1}-w^{*}\|_{H_{k}}^{2}
     +\|w_{k+1}-w_{k}\|_{H_{k}}^{2}
     +2\langle w_{k}-w_{k+1},\,H_{k}(w_{k+1}-w^{*})\rangle .
\tag{8}
\]

\noindent\textbf{[Step 5]}  Note that $w_{k}-w_{k+1}= \alpha H_{k}^{-1}g_{k}$, hence
\[
\langle w_{k}-w_{k+1},\,H_{k}(w_{k+1}-w^{*})\rangle
   =\alpha\langle g_{k},\,w_{k+1}-w^{*}\rangle .
\tag{9}
\]

\noindent\textbf{[Step 6]}  Plug \eqref{9} into \eqref{8} and rearrange:
\[
2\alpha\langle g_{k},\,w_{k+1}-w^{*}\rangle
   =\|w_{k}-w^{*}\|_{H_{k}}^{2}
    -\|w_{k+1}-w^{*}\|_{H_{k}}^{2}
    -\|w_{k+1}-w_{k}\|_{H_{k}}^{2}.
\tag{10}
\]

\noindent\textbf{[Step 7]}  Divide \eqref{10} by $2\alpha$ and substitute the result for the first term on the right‑hand side of \eqref{7}:
\[
f(w_{k})-f(w^{*})
   \le\frac{1}{2\alpha}\Bigl(
        \|w_{k}-w^{*}\|_{H_{k}}^{2}
       -\|w_{k+1}-w^{*}\|_{H_{k}}^{2}
       -\|w_{k+1}-w_{k}\|_{H_{k}}^{2}
        \Bigr)
      +\alpha\langle g_{k},H_{k}^{-1}g_{k}\rangle .
\tag{11}
\]

\noindent\textbf{[Step 8]}  The term $\|w_{k+1}-w_{k}\|_{H_{k}}^{2}$ can be written using the update:
\[
\|w_{k+1}-w_{k}\|_{H_{k}}^{2}
   =\|\alpha H_{k}^{-1}g_{k}\|_{H_{k}}^{2}
   =\alpha^{2}\langle g_{k},H_{k}^{-1}g_{k}\rangle .
\tag{12}
\]

\noindent\textbf{[Step 9]}  Insert \eqref{12} into \eqref{11} and cancel the identical $\alpha\langle g_{k},H_{k}^{-1}g_{k}\rangle$ terms:
\[
f(w_{k})-f(w^{*})
   \le\frac{1}{2\alpha}
        \Bigl(
          \|w_{k}-w^{*}\|_{H_{k}}^{2}
         -\|w_{k+1}-w^{*}\|_{H_{k}}^{2}
        \Bigr)
      -\frac{\alpha}{2}\langle g_{k},H_{k}^{-1}g_{k}\rangle .
\tag{13}
\]

\noindent\textbf{[Step 10]}  Sum \eqref{13} over $k=0,\dots,T-1$:
\[
\sum_{k=0}^{T-1}\!\bigl(f(w_{k})-f(w^{*})\bigr)
   \le\frac{1}{2\alpha}
        \Bigl(
          \|w_{0}-w^{*}\|_{H_{0}}^{2}
         -\|w_{T}-w^{*}\|_{H_{T}}^{2}
        \Bigr)
      -\frac{\alpha}{2}\sum_{k=0}^{T-1}\!\langle g_{k},H_{k}^{-1}g_{k}\rangle .
\tag{14}
\]

\noindent\textbf{[Step 11]}  Drop the non‑positive term $-\|w_{T}-w^{*}\|_{H_{T}}^{2}$ and use $H_{0}= \sqrt{\varepsilon}\,I$ to bound the first norm:
\[
\|w_{0}-w^{*}\|_{H_{0}}^{2}
   =\sum_{j=1}^{d}\frac{(w_{0,j}-w^{*}_{j})^{2}}{\sqrt{\varepsilon}}
   \le \frac{D^{2}}{\sqrt{\varepsilon}} .
\tag{15}
\]

\noindent\textbf{[Step 12]}  For the second sum, recall $H_{k}^{-1}=\operatorname{diag}\bigl(1/\sqrt{G_{k}+\varepsilon}\bigr)$. Hence
\[
\langle g_{k},H_{k}^{-1}g_{k}\rangle
   =\sum_{j=1}^{d}\frac{g_{k,j}^{2}}{\sqrt{G_{k,j}+\varepsilon}} .
\tag{16}
\]

\noindent\textbf{[Step 13]}  Since $G_{k,j}= \sum_{i=1}^{k}g_{i,j}^{2}$, the denominator is non‑decreasing in $k$, which yields
\[
\frac{g_{k,j}^{2}}{\sqrt{G_{k,j}+\varepsilon}}
   \le\frac{g_{k,j}^{2}}{\sqrt{G_{T,j}+\varepsilon}} .
\tag{17}
\]
Summing over $k$ gives
\[
\sum_{k=0}^{T-1}\frac{g_{k,j}^{2}}{\sqrt{G_{k,j}+\varepsilon}}
   \le\frac{\sum_{k=0}^{T-1}g_{k,j}^{2}}{\sqrt{G_{T,j}+\varepsilon}}
   =\frac{G_{T,j}}{\sqrt{G_{T,j}+\varepsilon}}
   \le\sqrt{G_{T,j}+\varepsilon} .
\tag{18}
\]

\noindent\textbf{[Step 14]}  Combining \eqref{14}–\eqref{18} and dividing by $T$ we obtain
\[
\frac{1}{T}\sum_{k=0}^{T-1}\!\bigl(f(w_{k})-f(w^{*})\bigr)
   \le
   \frac{D^{2}}{2\alpha T\sqrt{\varepsilon}}
   +\frac{\alpha}{2T}\sum_{j=1}^{d}\sqrt{G_{T,j}+\varepsilon}.
\tag{19}
\]

\noindent\textbf{[Step 15]}  Using the trivial bound $G_{T,j}\le T\,G_{\max}^{2}$ (from $\|g_{k}\|\le G_{\max}$) we have
\[
\sqrt{G_{T,j}+\varepsilon}\le\sqrt{T}\,G_{\max},
\qquad
\frac{1}{\sqrt{G_{T,j}+\varepsilon}}\le\frac{1}{\sqrt{\varepsilon}} .
\tag{20}
\]

\noindent\textbf{[Step 16]}  Plugging \eqref{20} into \eqref{19} yields the simplified $O(1/\sqrt{T})$ bound
\[
\frac{1}{T}\sum_{k=0}^{T-1}\!\bigl(f(w_{k})-f(w^{*})\bigr)
   \le
   \frac{D^{2}}{2\alpha T\sqrt{\varepsilon}}
   +\frac{\alpha G_{\max}\sqrt{d}}{2\sqrt{T}} .
\tag{21}
\]
Choosing $\alpha = D G_{\max}\sqrt{d}$ (or any constant order) balances the two terms and gives the clean rate
\[
\frac{1}{T}\sum_{k=0}^{T-1}\!\bigl(f(w_{k})-f(w^{*})\bigr)
   \le \frac{D G_{\max}\sqrt{d}}{\sqrt{T}} .
\tag{22}
\]
This completes the proof of Theorem~\ref{thm:adagrad}. \hfill $\square$

\section*{Rate Derivation (With Constants)}
From \eqref{19} we identify the two constants:
\[
C_{1}= \frac{D^{2}}{2\sqrt{\varepsilon}},\qquad
C_{2}= \frac{1}{2}\sum_{j=1}^{d}\sqrt{G_{T,j}+\varepsilon}.
\]
Thus the average suboptimality satisfies
\[
\frac{1}{T}\sum_{k=0}^{T-1}\bigl(f(w_{k})-f(w^{*})\bigr)
   \le \frac{C_{1}}{\alpha T}+\frac{\alpha C_{2}}{T}.
\]
Optimizing over $\alpha$ gives $\alpha^{*}= \sqrt{C_{1}/C_{2}}$ and the minimal bound
\[
\frac{2\sqrt{C_{1}C_{2}}}{T}
   = O\!\bigl(T^{-1/2}\bigr).
\]

\section*{Special Cases}
\begin{itemize}
\item \textbf{One‑dimensional case ($d=1$).} The bound reduces to
\[
\frac{1}{T}\sum_{k=0}^{T-1}\bigl(f(w_{k})-f(w^{*})\bigr)
   \le \frac{D G_{\max}}{\sqrt{T}} .
\]
\item \textbf{Sparse gradients.} If only $s\ll d$ coordinates receive non‑zero gradients, the sum $\sum_{j=1}^{d}\sqrt{G_{T,j}}$ scales with $s$, yielding a tighter $O(\sqrt{s/T})$ rate.
\item \textbf{Deterministic gradients.} When $g_{k}=\nabla f(w_{k})$ exactly, the same proof holds without expectation; the result is a deterministic regret bound.
\end{itemize}

\end{document}